% =====================================================================
%  第 11 章：强子结构
%  Gattringer & Lang, "Quantum Chromodynamics on the Lattice" (LNP 788)
%  由原始 pdftotext 转储（extract/ch11.txt）转换而来。
% =====================================================================
\chapter{强子结构}

在本节中，我们讨论一类重要的非能谱计算。强子结构通过合适的算符在强子态之间或强子态与真空之间的矩阵元来研究。这类性质中最简单的当属强子传播子的系数，它们与强子态和真空之间的矩阵元相联系，即所谓的衰变常数。它们描述强子的弱衰变性质。我们将较为详细地讨论衰变常数的格点计算，以此作为许多技术上有相似之处的问题的代表。
矢量流或轴矢量流在单强子态之间的矩阵元给出电磁形状因子和弱形状因子。进一步的矩阵元提供了关于半轻子衰变、夸克和胶子结构函数，以及诸如算符乘积展开等有效描述所需的其他信息。
所有这些都涉及场量，因而必须通过重整化方案与连续统量建立联系。格点微扰论提供了这样的关系。然而，其结果通常不适用于我们通常工作的、格点间距适中偏小的非微扰区域。因此，人们还依赖非微扰地确定重整化常数，这些常数使我们能与诸如 $\overline{\mathrm{MS}}$ 等连续统重整化方案中的可观测量建立联系。
在本章中，我们将对格点计算中获得矩阵元和重整化因子的相关方法作一个入门性的概述。

\section{低能参数}
在强子理论的早期，人们提出了若干假设来描述所观测到的现象。流代数假设将（味道 SU(3)）矢量流与轴矢量流联系在一组流的对易关系式（三元组）中。由于矢量流通过与电荷的关系而被归一化，这些对易关系也导致了轴矢量流的归一化。与弱相互作用耦合后，这一流描述了$\pi$介子的衰变等特征。部分守恒的轴矢流（PCAC）假设随后将轴矢量流的散度与$\pi$介子场联系起来，从而把$\pi$介子衰变常数定义为一个比例因子。该关系中的另一个因子是$\pi$介子质量的平方，这表明对于严格无质量的$\pi$介子，即戈德斯通玻色子，轴荷将守恒。根据这些假设，强子物理中的许多实验结果都可以得到解释。此后 QCD 被证明是其底层理论，原本的假设如今成为该理论的性质。在格点方法中，我们希望验证这些性质并测量相应的物理参数。
在本节中，我们只考虑 $u$ 和 $d$ 夸克。它们足够轻，使得手征味道 SU(2) 的显式破缺很小。

\subsection{算符定义}
在矩阵元的计算中，必须明确场与所观测到的粒子态之间的关系。在量子场论中，$\pi$介子场与所有具有$\pi$介子量子数 $J^{PC} = 0^{-+}$ 的态耦合。赝标量插入算符
\begin{equation}
P^a = \half\,\bar{\psi}\,\gamma_5\tau^a\psi
\end{equation}
具有这些量子数。这里我们把费米子场写为味道二重态 $\psi = (u, d)$，而 $\tau^a$，$a = 1, 2, 3$，是泡利矩阵。\footnote{在我们的讨论中，我们只考虑味道 SU(2) 对称群。}带电夸克双线性算符通过下式定义
\begin{equation}
\begin{aligned}
P^+ &\equiv P^1 - i P^2 = \half\,\bar{\psi}\gamma_5(\tau^1 - i\tau^2)\psi = \bar{d}\gamma_5 u , \\
P^- &\equiv P^1 + i P^2 = \half\,\bar{\psi}\gamma_5(\tau^1 + i\tau^2)\psi = \bar{u}\gamma_5 d .
\end{aligned}
\end{equation}

类似地，我们引入矢量场和轴矢量场
\begin{equation}
\begin{aligned}
V_\mu &= \half\,\bar{\psi}\gamma_\mu\psi , &\qquad
V_\mu^a &= \half\,\bar{\psi}\gamma_\mu\tau^a\psi , &\qquad
V_\mu^\pm &= V_\mu^1 \mp i V_\mu^2 , \\
A_\mu &= \half\,\bar{\psi}\gamma_\mu\gamma_5\psi , &\qquad
A_\mu^a &= \half\,\bar{\psi}\gamma_\mu\gamma_5\tau^a\psi , &\qquad
A_\mu^\pm &= A_\mu^1 \mp i A_\mu^2 .
\end{aligned}
\end{equation}
轴矢量流通过流代数中的对易关系与矢量流相联系。对于守恒矢量流，$A_\mu$ 和 $V_\mu$ 的归一化因此都被确定。
也可以直接在欧几里得空间中，通过场算符乘积期望值的积分恒等式得到等价的关系 [1]。在我们的表述中，我们将连续统量子场论的闵可夫斯基结果通过以 $t$ 替换 $it$ 转译为欧几里得形式。为了方便在连续统记法与格点记法之间切换，我们引入缩写 $\omega$ 表示空间傅里叶变换的归一化因子，即
\begin{equation}
\omega \equiv
\begin{cases}
\dfrac{1}{(2\pi)^3} & \text{for the continuum}, \\[2mm]
|\Lambda^3| & \text{for the lattice},
\end{cases}
\end{equation}
其中 $|\Lambda^3| = N^3$ 表示空间格点体积。
量子场论计算的结果在与实验数据比较之前需要重整化（参见第~\ref{sec:renorm} 节）。我们用上标 $(r)$ 表示按照连续统重整化方案重整化后的量。作为参考方案，我们采用标度 $\mu = 2$~GeV 处的 $\overline{\mathrm{MS}}$ 方案。我们定义重整化因子 $Z_P$、$Z_V$ 和 $Z_A$ 如下
\begin{equation}
P^{(r)a} = Z_P P^a , \qquad
V^{(r)a} = Z_V V^a , \qquad
A^{(r)a} = Z_A A^a
\end{equation}
并指出它们可能依赖于诸如格点作用量和标度等细节。在第~\ref{sec:renorm} 节中，我们将讨论如何确定联系裸格点场与重整化场的重整化因子。
夸克质量与凝聚项一样以乘法方式重整化
\begin{equation}
m^{(r)} = Z_m m , \qquad
\ev{\bar{\psi}\psi}^{(r)} = Z_S \ev{\bar{\psi}\psi} ,
\end{equation}
这引入了进一步的重整化因子 $Z_m$ 和 $Z_S$。由于组合 $m\bar{\psi}\psi$ 是重整化不变量，它们通过 $Z_S = 1/Z_m$ 相联系。若手征对称性成立，还有附加关系 $Z_S = Z_P$ 和 $Z_A = Z_V$。
如果费米子格点作用量不满足手征对称性，未重整化的裸夸克质量可能有一个加性贡献，称为残余质量，即 $m = m_{\mathrm{bare}} + m_{\mathrm{res}}$。
量子场论中的自由玻色子场满足标准的归一化。让我们假设渐近的、物理的$\pi$介子由这样一个分量 $\phi(x)^{(r)a}$ 的重整化三重态场描述。那么在欧几里得时空，与静止单$\pi$介子态 $\ket{\pi^b(p = 0)}$ 的矩阵元为
\begin{equation}
\braket{0}{\phi(x)^{(r)a}}{\pi^b(p = 0)} = \delta^{ab}\, e^{-M_\pi t} ,
\end{equation}
其中 $x = (\mathbf{x}, t)$。该场的传播子为
\begin{equation}
\braket{0}{\phi^{(r)a}(p = 0, t)\,\phi^{(r)b}(0)}{0} =
\frac{1}{2M_\pi\omega}\,\delta^{ab}\, e^{-M_\pi t} ,
\end{equation}
其中 $(p = 0, t)$ 表示在时间切片 $t$ 上投影到零空间动量（参见第 6.1.4 节）。然而，我们还需要把这个场与赝标量格点场算符 $P$ 联系起来。这将在下一节中完成。
在电弱相互作用的标准模型中，强子（如$\pi$介子或 K 介子）衰变到轻子的过程由一个——与强子质量相比——非常重的中间玻色子所介导。有效相互作用于是由一个将左手弱流与强子流耦合的拉格朗日量描述，例如对于衰变 $\pi \to e\,\nu_e$，有耦合项
\begin{equation}
\frac{G_F}{\sqrt{2}}\,\cos\theta_c\, \bigl(\bar{u}\gamma_\mu(1-\gamma_5)d\bigr)\,
\bigl(\bar{e}\gamma_\mu(1-\gamma_5)\nu_e\bigr) .
\end{equation}
这是弱哈密顿量一般算符乘积展开的一个例子，将在第~\ref{sec:weak} 节中更详细地讨论。
根据 (11.9)，$\pi$介子衰变给出轴矢量流在$\pi$介子态与真空之间的矩阵元。与物理的（重整化的）同位旋矢量$\pi$介子场的关系定义了$\pi$介子衰变常数 $F_\pi$，即
\begin{equation}
\partial_\mu A_\mu^{(r)a} = M_\pi^2 F_\pi\, \phi^{(r)a} .
\end{equation}
（也存在其他约定，例如相差一个因子 2。我们的定义对应于实验值 $F_\pi = 92.4(3)$~MeV [2]。）轴矢量流的散度充当$\pi$介子的插入场算符。这是闵可夫斯基空间中的一个算符恒等式；在欧几里得空间中，它作为期望值成立。
可以导出期望值的各种等价表达式。利用 (11.7) 我们得到所谓部分守恒的轴矢量流或 PCAC 关系：
\begin{equation}
\partial_\mu \braket{0}{A_\mu^{(r)a}(x)}{\pi^b(p = 0)} = \delta^{ab}\, M_\pi^2 F_\pi\, e^{-M_\pi t} .
\end{equation}
这些关系使我们能够从算符 $A_4$ 与 $\partial_t A_4$ 之间的成对关联函数导出衰变常数 $F_\pi$。例如，$\ev{A_4 A_4}$ 关联函数在大 $t$ 时由单$\pi$介子主导，因此具有渐近行为
\begin{equation}
\ev{A_4^{(r)+}(p = 0, t)\, A_4^{(r)-}(0)} \sim \frac{M_\pi F_\pi^2}{\sqrt{\omega}}\, e^{-M_\pi t} ,
\end{equation}
其中符号 $\sim$ 表示大 $t$ 时的渐近行为。K 介子衰变常数的计算与之类似，只需将其中的一个夸克替换为奇异夸克。更多细节将在第~\ref{sec:lowenergy} 节中讨论。
衰变常数只是可以获得的量之一。其他低能参数包括夸克质量和凝聚项。为此，我们在下一节中引入一个新的工具：Ward–Takahashi 恒等式。

\subsection{Ward 恒等式}
在配分函数的路径积分表达式中
\begin{equation}
Z = \int \mathcal{D}[\bar{\psi},\psi,U]\, e^{-S[\bar{\psi},\psi,U]}
  = \int \mathcal{D}[\bar{\psi}',\psi',U']\, e^{-S[\bar{\psi}',\psi',U']} ,
\end{equation}
场变量 $\bar{\psi},\psi, U$ 被积分。在第二步中，我们对路径积分使用了变换后的变量 $\bar{\psi}',\psi', U'$，这通常给出相同的结果。例外是影响积分测度的反常变换。这种利用变换后变量比较路径积分的策略已在第 7 章和第 9 章中成功应用。Ward 恒等式表达了配分函数（以及期望值）在这样一种场变量变换下的不变性。
让我们研究一个局域幺正变换的效果。任意插入算符 $O$ 的期望值为
\begin{equation}
\braket{0}{O}{0} =
\frac{1}{Z}\int \mathcal{D}[\bar{\psi},\psi,U]\, O[\bar{\psi},\psi,U]\, e^{-S[\bar{\psi},\psi,U]} .
\end{equation}
我们考虑此路径积分中费米子场的无穷小对称变换 $\psi \to \psi + \delta\psi$，$\bar{\psi} \to \bar{\psi} + \delta\bar{\psi}$。结果必须在这样一种积分变量变换下不变。对于非反常变换，积分测度不变，于是得到形如下的（Ward）恒等式
\begin{equation}
0 = \braket{0}{\delta O}{0} - \braket{0}{O\,\delta S}{0} ,
\end{equation}
其中 $\delta O$ 和 $\delta S$ 分别表示算符 $O$ 和作用量 $S$ 在变换下的线性变化。最简单的情形 $O = 1$ 为
\begin{equation}
\braket{0}{\delta S}{0} = 0 ,
\end{equation}
它给出与经典诺特定理守恒律类似的关系。对于某些变换，泛函积分测度并非不变。这便产生所谓的反常贡献（参见第 7 章和第 9 章）。
让我们首先在熟悉的连续统形式下、对 $O = 1$ 导出这些关系。稍后我们将指出与格点记法的差别。我们研究作用量（$M = \mathrm{diag}(m_u, m_d)$ 表示味道空间中的质量矩阵）的变换
\begin{equation}
S = \int d^4x\, \bar{\psi}(\gamma_\mu D_\mu + M)\psi , \qquad
D_\mu = \partial_\mu + iA_\mu ,
\end{equation}
考虑无穷小 SU(2) 味道变换，对应于有限变换 (7.13)、(7.14)、(7.15) 和 (7.16)。线性化的无穷小变换可写为
\begin{equation}
\begin{aligned}
\psi(x) &\to \psi'(x) = \bigl(1 + i\varepsilon(x)\lambda\bigr)\psi(x) , \\
\bar{\psi}(x) &\to \bar{\psi}'(x) = \bar{\psi}(x)\bigl(1 + i\varepsilon(x)\hat{\lambda}\bigr) ,
\end{aligned}
\end{equation}
其中 $\lambda$ 和 $\hat{\lambda}$ 是 Dirac 空间与味道空间中矩阵的乘积，例如 $\lambda = 1$、$\tau^a$、$\gamma_5$、$\gamma_5\tau^a$，而 $\hat{\lambda} = -1$、$-\tau^a$、$\gamma_5$、$\gamma_5\tau^a$。我们假设 $\varepsilon(x)$ 是一个在某有界区域外光滑地趋于零的函数。这一点很重要，因为我们使用了分部积分，边界项不应有贡献。
到 $O(\varepsilon)$ 阶，作用量 (11.17) 在这样一种无穷小变换下的变化为
\begin{equation}
\delta S = i \int d^4x\,
\bar{\psi}\Bigl[\varepsilon\hat{\lambda}\gamma_\mu\partial_\mu + \gamma_\mu\lambda\,\partial_\mu\varepsilon
+ i\varepsilon A_\mu(\hat{\lambda}\gamma_\mu + \gamma_\mu\lambda)
+ \varepsilon(\hat{\lambda}M + M\lambda)\Bigr]\psi .
\end{equation}
这里我们只研究具有以下性质的变换
\begin{equation}
\hat{\lambda}\gamma_\mu + \gamma_\mu\lambda = 0 ,
\end{equation}
由于规范场对味道和 Dirac 结构都不敏感，我们去掉了含 $A_\mu$ 的项。我们还有
\begin{equation}
\partial_\mu(\varepsilon\psi) = (\partial_\mu\varepsilon)\psi + \varepsilon\,\partial_\mu\psi ,
\end{equation}
由此得到
\begin{equation}
\begin{aligned}
\delta S &= i \int d^4x \Bigl[\bigl(\partial_\mu\varepsilon(x)\bigr)\bar{\psi}\gamma_\mu\lambda\psi
+ \varepsilon(x)\bar{\psi}\bigl(\hat{\lambda}M + M\lambda\bigr)\psi\Bigr] \\
&= i \int d^4x\, \varepsilon(x)\Bigl[-\partial_\mu\bigl(\bar{\psi}\gamma_\mu\lambda\psi\bigr)
+ \bar{\psi}\bigl(\hat{\lambda}M + M\lambda\bigr)\psi\Bigr] .
\end{aligned}
\end{equation}
在最后一步中我们使用了分部积分，假设边界项没有贡献，正如我们选择 $\varepsilon(x)$ 所保证的那样。由于这一选择是任意的，条件 (11.16) 必须对任意 $x$ 成立。由此得到算符恒等式
\begin{equation}
\partial_\mu\bigl(\bar{\psi}\gamma_\mu\lambda\psi\bigr) = \bar{\psi}\bigl(\hat{\lambda}M + M\lambda\bigr)\psi .
\end{equation}
然后，根据 $\lambda$ 和 $\hat{\lambda}$ 的选择（但须满足 (11.20)），我们得到关系
\begin{equation}
\lambda = 1,\ \hat{\lambda} = -1 \quad\Rightarrow\quad
\partial_\mu\bigl(\bar{\psi}\gamma_\mu\psi\bigr) = 0 ,
\end{equation}
\begin{equation}
\lambda = \tau^a,\ \hat{\lambda} = -\tau^a \quad\Rightarrow\quad
\partial_\mu\bigl(\bar{\psi}\gamma_\mu\tau^a\psi\bigr) = \bar{\psi}\,[M,\tau^a]\,\psi ,
\end{equation}
\begin{equation}
\lambda = \hat{\lambda} = \gamma_5 \quad\Rightarrow\quad
\partial_\mu\bigl(\bar{\psi}\gamma_\mu\gamma_5\psi\bigr) = 2\,\bar{\psi}M\gamma_5\psi
\;(+\mathrm{anomaly}) ,
\end{equation}
\begin{equation}
\lambda = \hat{\lambda} = \gamma_5\tau^a \quad\Rightarrow\quad
\partial_\mu\bigl(\bar{\psi}\gamma_\mu\gamma_5\tau^a\psi\bigr)
= \bar{\psi}\,\{M,\tau^a\}\,\gamma_5\psi ,
\end{equation}
其中我们用到质量矩阵 $M$ 与 $\gamma_5$ 对易：$M\gamma_5 = \gamma_5 M$。
利用 (11.1) 和 (11.3) 对味道单态和非单态矢量流、轴矢量流以及赝标量场的简写记号，上述关系便具有如下形式
\begin{equation}
\partial_\mu V_\mu = 0 ,
\end{equation}
\begin{equation}
\partial_\mu V_\mu^a = \half\,\bar{\psi}\,[M,\tau^a]\,\psi ,
\end{equation}
\begin{equation}
\partial_\mu A_\mu = \bar{\psi}M\gamma_5\psi \;(+\mathrm{anomaly}) ,
\end{equation}
\begin{equation}
\partial_\mu A_\mu^a = \half\,\bar{\psi}\,\{M,\tau^a\}\,\gamma_5\psi .
\end{equation}

关于味道单态矢量流的方程 (11.28) 是总重子数守恒律。在另外三个关系中，我们发现诺特对称性可能被质量项破坏。对于质量简并的夸克，存在同位旋对称性，因此 $[M, \tau^a] = 0$。于是 (11.29) 定义了味道矢量流的守恒。在 (11.30) 中，当为量子化理论写出时，由于泛函积分中积分测度的非不变性，必须加上反常项（参见第 7.1、7.3 和 9.4 节）。方程 (11.31) 被称为非单态轴矢 Ward 恒等式（AWI）：味道非单态轴矢量流的散度与赝标量场相联系。
我们再次强调，由于 (11.16)，这些关系作为期望值成立。由于 $O$ 的选择是任意的，我们可以选择算符在某个紧致的、任意小的时空区域内非零且为常数，而在区域外为零。因此，(11.28)、(11.29)、(11.30) 和 (11.31) 可被视为局域地成立。这些关系在量子化和重整化下被保持（除反常之外），因此对重整化量同样成立。
对于简并的费米子质量 $m$，AWI (11.31) 具有如下形式
\begin{equation}
\partial_\mu A_\mu^{(r)a} = 2\,m^{(r)} P^{(r)a} .
\end{equation}
这一关系在许多方面都很有用。将其与 (11.10) 比较并在真空与$\pi$介子之间求值，我们发现
\begin{equation}
M_\pi^2 F_\pi \braket{0}{\phi^{(r)a}}{\pi}
= \braket{0}{\partial_\mu A_\mu^{(r)a}}{\pi}
= 2\,m^{(r)} \braket{0}{P^{(r)a}}{\pi} .
\end{equation}
这使我们能够从 $P^{(r)a}$ 和 $A_\mu^{(r)a}$ 关联函数的渐近行为确定重整化夸克质量 $m^{(r)}$ 和$\pi$介子衰变常数（参见第~\ref{sec:lowenergy} 节）。
进一步利用 (11.32) 会得到另一个有趣的关系。我们利用 (11.22)，将变换 (11.31) 乘以一个局域赝标量算符，并对其在格拉斯曼场上取期望值。具体地，我们使用 (11.18)，其中
\begin{equation}
\begin{aligned}
\lambda = \hat{\lambda} &= \gamma_5\,\half(\tau^1 + i\tau^2) , \\
O(0) &= \bar{\psi}(0)\gamma_5\,\half(\tau^1 - i\tau^2)\psi(0)
= \bar{d}(0)\gamma_5 u(0) = P^+(0) .
\end{aligned}
\end{equation}
这个算符正是赝标量场。在手征极限 $M = 0$ 下，方程 (11.15) 结合 (11.22) 现在给出
\begin{equation}
-\int d^4x\, \ev{\partial_\mu A_\mu^{(r)-}(x)\, P^+(0)}
= \ev{\bar{u}(0)u(0) + \bar{d}(0)d(0)} .
\end{equation}
\begingroup\emergencystretch=4em 在左边插入一组完备态，两点函数被$\pi$介子饱和，表达式正比于\linebreak[0] $\ev{P^{-}(x)P^{+}(0)}$。左边变为\par\endgroup
\begin{equation}
-\frac{F_\pi^2 M_\pi^4}{2\,m^{(r)}}\int d^4x\, \phi^-(x)\phi^+(0)
= -\frac{F_\pi^2 M_\pi^2}{m^{(r)}} ,
\end{equation}
其中我们用到该积分是在零动量处动量空间传播子 $(p^2 + M_\pi^2)^{-1}$ 的两倍。
(11.35) 的右边是凝聚项 $\Sigma^{(r)}$
\begin{equation}
\ev{\bar{u}(0)u(0) + \bar{d}(0)d(0)} \equiv N_f \Sigma^{(r)} ,
\end{equation}
在我们的情形中 $N_f = 2$。
我们最终得到一个对重整化夸克质量和重整化凝聚项 $\Sigma^{(r)}$ 成立的关系：
\begin{equation}
F_\pi^2 M_\pi^2 = -m^{(r)}\, N_f\, \Sigma^{(r)} .
\end{equation}
$\pi$介子衰变常数 $F_\pi$ 和$\pi$介子质量 $M_\pi$ 不依赖于重整化方案；它们是物理的、可测量的量。在我们的归一化下，$F_\pi$ 的实验值接近 93~MeV。
方程 (11.38) 因开创性的论文 [3] 而被称为 Gell–Mann–Oakes–Renner (GMOR) 关系，该文讨论了由夸克质量项引起的对称性破缺机制。注意该方程可能有 $O(m^2)$ 阶的修正项。这是一个系统性展开的最低阶，该展开导致一种称为手征微扰论 [4, 5] 的有效场论。
当人们试图推导连续统恒等式的格点类比 [6, 7] 时，会遇到一个问题。对于不满足手征对称性的作用量，如威尔逊作用量，格点作用量的手征变分不能仅用轴矢流和轴矢密度来表达 [6]，还会出现额外的项。不过，人们随后论证，重整化格点场的关联函数应当收敛到相应的连续统关联函数。这使得人们可以从诸如 (11.58) 这样的、针对非零传播距离的期望值中提取信息。

\subsection{格点上的朴素流与守恒流}
\label{sec:naivecurrents}
现在我们转向格点表述，考虑构造与形如 (11.18) 的对称变换相对应的格点流，即
\begin{equation}
\delta\psi(n) = i\varepsilon(n)\lambda\psi(n) , \qquad
\delta\bar{\psi}(n) = i\bar{\psi}(n)\varepsilon(n)\hat{\lambda} .
\end{equation}
在此变换下，作用量
\begin{equation}
S = \sum_{n,m} \bar{\psi}(n)\, D(U; n|m)\, \psi(m)
\end{equation}
发生变化，其变化的线性项为
\begin{equation}
\delta S = i \sum_{n,m}
\Bigl[\bar{\psi}(n)\, D(U; n|m)\, \varepsilon(m)\lambda\, \psi(m)
+ \varepsilon(n)\,\bar{\psi}(n)\hat{\lambda}\, D(U; n|m)\, \psi(m)\Bigr] .
\end{equation}

让我们首先考虑威尔逊费米子 (5.51)，取 $\lambda = \tau^a$ 和 $\hat{\lambda} = -\tau^a$。威尔逊作用量变换中的 $O(\varepsilon)$ 项必须对任意 $\varepsilon(n)$ 消失。在对格点上的求和进行指标平移后，这可以化为局域形式。我们最终得到与 (11.29) 等价的方程，其中连续统导数被最近邻差分所取代
\begin{equation}
\sum_{\mu=1}^{4} \frac{1}{a}\Bigl[V_\mu^a(n + \hat{\mu}) - V_\mu^a(n)\Bigr]
= \half\,\bar{\psi}(n)\,[M, \tau^a]\,\psi(n) ,
\end{equation}
其中格点矢量流为
\begin{equation}
V_\mu^a(n) = \half \Bigl[
\bar{\psi}(n + \hat{\mu})(1 + \gamma_\mu)U_\mu(n)^\dagger\tau^a\psi(n)
- \bar{\psi}(n)(1 - \gamma_\mu)U_\mu(n)\tau^a\psi(n + \hat{\mu})
\Bigr] .
\end{equation}

我们可以将 (11.42) 写为
\begin{equation}
\Delta_\mu V_\mu^a(n) = \half\,\bar{\psi}(n)\,[M, \tau^a]\,\psi(n) ,
\end{equation}
其中 $\Delta_\mu$ 是格点前向导数 $\Delta_\mu f(n) \equiv \bigl(f(n + \hat{\mu}) - f(n)\bigr)/a$。这与相应的连续统表达式 (11.29) 一致（相差格点间距的高阶项）。对于质量相同的 $u$ 和 $d$ 夸克，该流与连续统情形一样守恒。
然而，对于破坏手征对称性的费米子作用量（如威尔逊作用量），轴矢变换的这种简单等价关系不再成立。存在随 $1/a$ 发散的对称性破缺项，因此不可能在夸克质量为零时定义一个精确守恒的轴矢量流。
在第 7 章中，我们讨论了满足 Ginsparg–Wilson 关系从而在格点上具有手征对称性的 Dirac 算符。相应的算符，例如 overlap 算符，要复杂得多。因此，考虑在未明确指定具体形式的情况下为一般作用量构造诺特流的问题是有价值的。为此，我们遵循 [8]，作为一个例子讨论形如 (11.18) 和 (11.24) 的、$\lambda = 1$、$\hat{\lambda} = -1$ 的简单 $n$ 依赖标量变换。
变分 (11.41) 具有如下形式（这里 $k$ 是位置指标）
\begin{equation}
\begin{aligned}
\delta S &= i a \sum_{k,\mu} \bigl(\Delta_\mu\varepsilon(k)\bigr) J_\mu(k)
= i \sum_{k,\mu} \bigl(\varepsilon(k + \hat{\mu}) - \varepsilon(k)\bigr) J_\mu(k) \\
&= i \sum_{k,\mu} \varepsilon(k) \bigl(J_\mu(k - \hat{\mu}) - J_\mu(k)\bigr)
= -i a \sum_{k,\mu} \varepsilon(k)\, \Delta_\mu^* J_\mu(k) ,
\end{aligned}
\end{equation}
其中 $\Delta_\mu$ 和 $\Delta_\mu^*$ 分别表示格点前向和后向导数。如果作用量在这样一种变量变换下不变，即 $\delta S = 0$，我们就可以将 $J_\mu(k)$ 等同于格点上的守恒诺特流。

为此，让我们考虑对每一个单独的链（link）变量的变换：
\begin{equation}
U_\mu'(n) = U_\mu^{(\alpha)}(n) \equiv e^{i\alpha_\mu(n)} U_\mu(n)
\equiv e^{i\varepsilon(n)} U_\mu(n)\, e^{-i\varepsilon(n + \hat{\mu})} ,
\end{equation}
其中 $\alpha_\mu(n) = -\Delta_\mu\varepsilon(n)$。我们改写这一变换是为了利用 Dirac 算符在规范变换下的变换性质，于是
\begin{equation}
\begin{aligned}
&U_\mu(n) \to U_\mu'(n) \qquad\Rightarrow\qquad\\
&D(U; n|m) \to D(U'; n|m) = e^{i\varepsilon(n)}\, D(U; n|m)\, e^{-i\varepsilon(m)} .
\end{aligned}
\end{equation}

取变换后与变换前的 Dirac 算符之差，在 $O(\varepsilon)$ 阶我们得到：
\begin{equation}
D(U'; n|m) - D(U; n|m)
= i\bigl(D(U; n|m)\,\varepsilon(m) - \varepsilon(n)\, D(U; n|m)\bigr) ,
\end{equation}
其中我们只考虑了 $\varepsilon$ 的最低阶。而这正是 (11.41) 的核。另一种方式，我们也可以通过对 $\alpha_\mu$ 求导来确定这一最低阶之差：
\begin{equation}
\begin{aligned}
D(U'; n|m) - D(U; n|m)
&= -\sum_{k,\mu} \alpha_\mu(k)
\left.\left(\frac{\partial D\bigl(U_\mu^{(\alpha)}(k); n|m\bigr)}
{\partial\alpha_\mu(k)}\right)\right|_{\alpha=0} \\
&= \sum_{k,\mu} \bigl(\varepsilon(k + \hat{\mu}) - \varepsilon(k)\bigr)
\left.\left(\frac{\partial D\bigl(U_\mu^{(\alpha)}(k); n|m\bigr)}
{\partial\alpha_\mu(k)}\right)\right|_{\alpha=0} .
\end{aligned}
\end{equation}

与 (11.45) 比较，可以确定诺特流为
\begin{equation}
\begin{aligned}
J_\mu(k) &= \sum_{n,m} \bar{\psi}(n)\, K_\mu(k; n|m)\, \psi(m) , \\
K_\mu(k; n|m) &\equiv -i
\left.\left(\frac{\partial D\bigl(U_\mu^{(\alpha)}(k); n|m\bigr)}
{\partial\alpha_\mu(k)}\right)\right|_{\alpha=0} .
\end{aligned}
\end{equation}
我们发现，如果把 Dirac 算符想象为连接 $n$ 与 $m$ 的路径之和，那么所有经过链 $U_\mu(k)$ 的路径都对 $J_\mu(k)$ 有贡献。从构造来看，$J_\mu(k)$ 的规范不变性是显然的。在实际应用中，导数必须通过取差分在数值上实现。将此构造应用于威尔逊 Dirac 算符，我们便恢复出如 (11.43) 中那样的简单局域流。
对于手征对称的作用量，可以构造守恒且协变变换的矢量流以及协变的标量密度。它们随后满足与连续统中类似的 Ward 恒等式。
遵循我们例子中概述的策略，在 [8, 9] 中，人们为一般的、精确手征不变的作用量（如 overlap 作用量）构造了手征协变且守恒的流。详细的推导我们请读者参阅原始论文，这里只引用结果。利用缩写
\begin{equation}
\hat{\gamma}_5 = \gamma_5 (1 - D) ,
\end{equation}
我们有
\begin{equation}
\begin{aligned}
V_\mu^a(k) &=
\frac{1}{4}\sum_{n,m} \bar{\psi}(n)\Bigl[K_\mu(k; n|m)
- \gamma_5\sum_j K_\mu(k; n|j)\,(\hat{\gamma}_5)_{jm}\Bigr]\tau^a\psi(m) , \\[2mm]
A_\mu^a(k) &=
\frac{1}{4}\sum_{n,m} \bar{\psi}(n)\Bigl[-\gamma_5 K_\mu(k; n|m)
+ \sum_j K_\mu(k; n|j)\,(\hat{\gamma}_5)_{jm}\Bigr]\tau^a\psi(m) , \\[2mm]
S^a(k) &=
\frac{1}{4}\sum_{n,m} \bar{\psi}(n)\Bigl[\delta_{nk}\bigl(1 - \half D\bigr)_{km}\\
&\qquad + \half\sum_j \bigl(1 - \half D\bigr)_{nj}\delta_{jk}\delta_{mk}
+ (\hat{\gamma}_5)_{jk}(\hat{\gamma}_5)_{km}\Bigr]\tau^a\psi(m) , \\[2mm]
P^a(k) &=
\frac{1}{4}\sum_{n,m} \bar{\psi}(n)\gamma_5\Bigl[\delta_{nk}\bigl(1 - \half D\bigr)_{km}\\
&\qquad + \half\sum_j \bigl(1 - \half D\bigr)_{nj}(\hat{\gamma}_5)_{jk}\delta_{mk}
+ \delta_{jk}(\hat{\gamma}_5)_{km}\Bigr]\tau^a\psi(m) .
\end{aligned}
\end{equation}
这些守恒流以及等价构造的标量、赝标量密度在手征变换下协变地变换，与连续统中一样。对于守恒矢量流，重整化是平凡的（电荷守恒，且重整化常数满足 $Z_A = Z_V = 1$）。一个重要的优点是，这些流和密度（维度-3 算符）自动是 $O(a)$ 改进的，并且分别不与维度-4 或维度-3 算符混合。
由于方案 (11.50) 的实现相当繁琐，人们使用更简单的表达式
\begin{equation}
\begin{aligned}
\bar{V}_\mu^a(k) &=
\half \sum_m \bar{\psi}(k)\gamma_\mu\bigl(1 - \half D\bigr)_{km}\tau^a\psi(m) , \\[1mm]
\bar{A}_\mu^a(k) &=
\half \sum_m \bar{\psi}(k)\gamma_\mu\gamma_5\bigl(1 - \half D\bigr)_{km}\tau^a\psi(m) , \\[1mm]
\bar{S}^a(k) &=
\half \sum_m \bar{\psi}(k)\bigl(1 - \half D\bigr)_{km}\tau^a\psi(m) , \\[1mm]
\bar{P}^a(k) &=
\half \sum_m \bar{\psi}(k)\gamma_5\bigl(1 - \half D\bigr)_{km}\tau^a\psi(m)
\end{aligned}
\end{equation}
\begingroup\emergencystretch=4em 它们也是协变的，但并非守恒，因此不受 Ward 恒等式的约束 [8]。\linebreak[1]基于 Ginsparg–Wilson 费米子的直接构造还可参见 [10, 11]。\par\endgroup
在距离大于 Dirac 算符作用范围的关联函数中，(11.52) 中的协变密度可以被更简单的、点状表达式所取代
\begin{equation}
\begin{aligned}
S^a(k) &\approx \half\, \frac{1}{1 - m/2}\,\bar{\psi}(k)\tau^a\psi(k) , \\
P^a(k) &\approx \half\, \frac{1}{(1 - m/2)(1 - m^2/4)}\,\bar{\psi}(k)\gamma_5\tau^a\psi(k) .
\end{aligned}
\end{equation}
这是因为传播子涉及有质量 Dirac 算符的逆 $D_m^{-1}$，而由 (7.41) 和 (7.42) 我们知道
\begin{equation}
1 - \half D = \frac{1}{1 - m/2}\Bigl(1 - \half D_m\Bigr) .
\end{equation}
因此，(11.53) 中的插入因子 $1 - \half D$ 可以用有质量 Dirac 算符 $D_m$ 来表达。在考虑插入算符 (11.53) 的关联函数时，这些插入因子会产生乘积 $D_m D_m^{-1}$，它们只在局域处有贡献（“接触项”），可以被省略 [8]。
在连续极限中，这些流以 $Z_V$ 和 $Z_A$ 之类的因子乘法地重整化，我们将在第~\ref{sec:renorm} 节中讨论这些因子的确定。

\subsection{由关联函数得到的低能参数}
\label{sec:lowenergy}
在上一节中，我们已经说明，关联函数的连续统恒等式可以被转译到格点上，尽管在格点上识别流和密度并非直截了当。我们可以预期，关联函数的渐近行为是可以比较的。因此，我们可以利用如 (11.12) 之类的方程，从某些关联函数及其比值的渐近（长时距）行为来确定低能常数。
重整化夸克质量可以利用轴矢 Ward 恒等式 (11.31)、(11.33)，并考虑下式的（大 $t$）渐近行为来确定
\begin{equation}
\frac{1}{2}\,
\frac{\ev{0}{\partial_t A_4^{-(r)}(p = 0, t)\,X(0)}{0}}
{\ev{0}{P^{-(r)}(p = 0, t)\,X(0)}{0}}
= \frac{1}{2}\,
\frac{Z_A \ev{0}{\partial_t A_4^-(p = 0, t)\,X(0)}{0}}
{Z_P \ev{0}{P^-(p = 0, t)\,X(0)}{0}}
\sim m^{(r)} .
\end{equation}
对于 $X$，人们选择一个与$\pi$介子耦合的插入算符，通常取 $X = P$。分子和分母在渐近区都主要由$\pi$介子中间态主导，渐近的 $t$ 依赖被消去。在计算中，一旦渐近行为建立起来，就可以识别出一个平台。在此关系中，需要知道将格点插入算符与连续统归一化联系起来的重整化常数 $Z_A$ 和 $Z_P$。它们可以单独确定，这一点将在后面讨论。
忽略重整化常数，人们常常只确定未重整化关联函数的渐近比值：
\begin{equation}
\frac{1}{2}\,
\frac{\ev{0}{\partial_t A_4^-(p = 0, t)\,X(0)}{0}}
{\ev{0}{P^-(p = 0, t)\,X(0)}{0}}
\sim m_{\mathrm{AWI}} .
\end{equation}
由于其来源，这个比值被称为 AWI 质量 $m_{\mathrm{AWI}}$ 或 PCAC 质量。我们已经在第 9.1.4 节的 Symanzik 改进方案中把 (11.57) 用作确定系数 $c_{sw}$ 和 $c_A$ 的工具。
格点作用量的夸克质量可能经历两种重整化。如果作用量不满足手征对称性，可能存在一种被称为残余质量 $m_{\mathrm{res}}$ 的加性重整化，且 $m_{\mathrm{AWI}} = m_{\mathrm{bare}} + m_{\mathrm{res}}$。于是我们有 $m^{(r)} = (Z_A / Z_P)\, m_{\mathrm{AWI}}$。例如，相对于连续统 $\overline{\mathrm{MS}}$ 方案，存在乘法重整化因子 $m^{(r)} = Z_m m$。对于手征对称的作用量，$Z_m Z_P = 1$，因此以格点单位计的裸质量为 $m = Z_A\, m_{\mathrm{AWI}}$。
AWI 质量在手征对称性成立的耦合常数空间中的那些点处为零。这在存在加性质量重整化的情形中特别有用，因为它可以识别出残余质量（见下文）。对于威尔逊作用量，人们借助它确定跳跃参数的临界值 $\kappa_c(\beta)$：它被定义为 $m_{\mathrm{AWI}}$ 为零且$\pi$介子无质量时的取值（也可参见第 9.1.4 节）。
$\pi$介子衰变常数可以例如从轴矢流关联函数 (11.12) 得到：
\begin{equation}
Z_A^2 \ev{A_4^+(p = 0, t)\, A_4^-(0)} \sim \frac{M_\pi F_\pi^2}{\sqrt{\omega}}\, e^{-M_\pi t} .
\end{equation}
由 (11.33)，我们还得到
\begin{equation}
\ev{0}{P^+(p = 0, t)\,P^-(0)}{0} \sim \frac{A}{\sqrt{\omega}}\, e^{-M_\pi t}
\qquad\text{with}\qquad
A = \frac{M_\pi^3 F_\pi^2}{4\,m^{(r)2} Z_P^2} .
\end{equation}
从渐近行为的系数我们得到
\begin{equation}
F_\pi = 2\,Z_P\, m^{(r)}\, \sqrt{\frac{A}{M_\pi^3}}
= 2\,Z_A\, m_{\mathrm{AWI}}\, \sqrt{\frac{A}{M_\pi^3}} .
\end{equation}
一旦 $F_\pi$、$M_\pi$ 和 $m^{(r)}$ 已知，GMOR 方程 (11.38) 就可以用来计算夸克凝聚项。人们也可以使用如下形式的关联函数
\begin{equation}
Z_A Z_P\, \bigl|\ev{0}{A_4^+(p = 0, t)\,P^-(0)}{0}\bigr| \sim \frac{|\Sigma^{(r)}|}{\sqrt{\omega}}\, e^{-M_\pi t}
\end{equation}
来确定它（更多此类关系参见 [12]）。

\section{重整化}
\label{sec:renorm}
\subsection{为什么需要重整化？}
在量子化一个场论时，必须对其进行正则化，而格点表述是一种可能的紫外正则子。将计算出的量与物理可观测量对应起来便完成了重整化过程。不同的正则子和作用量需要不同的重整化参数。一些可观测量，如强子质量，可以直接作为无量纲数 $a\,m_{\mathrm{phys}}$ 从关联函数的指数衰减中测量出来。将这样一个数与相应粒子的物理质量 $m_{\mathrm{phys}}$ 比较，就可以确定格点间距 $a$ 并建立标度（参见第 6.3 节）。其他量，如衰变常数或形状因子，涉及场，它们与实验及连续统归一化的关系更为复杂。
矩阵元以及诸如夸克质量或凝聚项之类的低能参数在 QCD 唯象学中起着重要作用。然而，它们的值依赖于某种重整化方案中的定义。由于人们希望与物理的、实验测量的量比较，就必须将这些量联系到量子化理论底层表述的参数。为此，必须确定可观测量随标度 $a$ 的标度性质以及它们在重整化下的行为。许多量本质上是标度依赖的，甚至当截断参数 $a$ 被移除时发散。我们可以区分以下情形：
\begin{itemize}
\item 有限算符，例如矢量流和轴矢量流，或标量密度与赝标量密度之比：乘法重整化因子 $Z_V$、$Z_A$ 以及比值 $Z_S/Z_P$ 在正则化参数 $a$ 趋于零时取有限值，因此它们应当与标度无关。对于守恒流和手征对称的作用量，$Z_V = Z_A = 1$。取决于 Dirac 算符，使用守恒流可能很复杂（参见第~\ref{sec:naivecurrents} 节），因此人们依赖于更简单的、例如点状的流，此时 $Z$ 因子不能先验地知道。
\item 对数发散的算符（例如标量和赝标量密度，或改变味道 $\Delta F = 2$ 的四费米子算符）是标度依赖的。
\item 幂次发散的算符：这些出现在具有内禀质量标度的正则化中，例如格点正则化。此时与更低维算符的混合成为可能，导致幂次发散 $\order{1/a^n}$ [13]。例子是在第 10.4 节中讨论的重夸克有效理论（HQET）中的 $\order{1/m}$ 项。
\end{itemize}

在正则化中，原始理论的一些对称性会丢失。在格点方法中，这些例如是连续的时空变换和（对某些作用量而言的）手征对称性。精确的 GW 费米子受到其手征对称性的保护，由它们构造的算符的重整化比例如威尔逊费米子的更简单。手征对称性蕴含重整化常数之间的若干关系，例如 $Z_A = Z_V$ 和 $Z_S = Z_P$。
为了在格点上非微扰地计算重整化常数，需要一个既能在格点蒙特卡罗模拟中实现、又能在连续统微扰论中实现的重整化方案。后一性质对于将格点结果转换到连续统方案是必要的。最受青睐的连续统方案是修正最小减除 $\overline{\mathrm{MS}}$ 方案。在这样一种与质量无关的方案中，重整化因子只依赖于一个归一化质量 $\mu$ 和耦合常数，而不依赖于夸克质量。

\subsection{用罗马–南安普顿方法重整化}
由于人们将格点结果与连续统结果比较，本可以利用格点微扰论（参见 [13] 及其中的参考文献）。这确已做过，但在许多情形下收敛性太差，无法得到可靠的结果。因此人们常常依赖非微扰方法 [14]。
其思想是将非微扰格点计算中确定的裸格点关联函数、耦合常数和质量与所谓正则化无关（RI）方案（通常也称为 RI/MOM 方案）中的量进行比较。RI 量与 $\overline{\mathrm{MS}}$ 方案中定义的量之间的联系随后可以在连续统微扰论中建立 [15]。
作为例子，我们讨论夸克双线性算符的重整化因子。更明确地，让我们考虑局域的、味道非单态夸克场双线性算符
\begin{equation}
O_\Gamma \equiv \bar{u}\,\Gamma\, d .
\end{equation}
这里 $\Gamma$ 表示一个 Clifford 代数矩阵。根据它们的洛伦兹对称性，我们用 $S$、$V$、$A$、$T$ 和 $P$ 表示五种类型的 $\Gamma$，分别对应标量、矢量、轴矢量、张量和赝标量（分别对应 $1$、$\gamma_\mu$、$\gamma_\mu\gamma_5$、$\frac{i}{2}[\gamma_\mu, \gamma_\nu]$ 和 $\gamma_5$）。
人们研究在特定动量值 $p^2 = \mu^2$ 处夸克场 $\ket{p}$ 之间双线性夸克算符的期望值 $\braket{p}{O_\Gamma}{p}$（即在零动量转移处），并将它们与相应的树级矩阵元匹配：
\begin{equation}
Z_\Gamma\, \braket{p}{O_\Gamma}{p}\big|_{p^2 = \mu^2} = \braket{p}{O_\Gamma}{p}_0\big|_{p^2 = \mu^2} .
\end{equation}
重整化常数 $Z_\Gamma$ 是相互作用情形与自由情形之间的比例因子。
预期这一程序在如下窗口内有效
\begin{equation}
\Lambda_{\mathrm{QCD}}^2 \ll \mu^2 \ll 1/a^2 ,
\end{equation}
其中离散化效应可以被忽略，因为重整化标度 $\mu$ 与格点截断 $1/a$ 相比很小。此时（少圈的）连续统微扰论可以用来连接不同的方案，因为 $\mu$ 远大于 QCD 标度参数 $\Lambda_{\mathrm{QCD}}$。与 $\overline{\mathrm{MS}}$ 方案比较时，典型值为 $\mu = 2$~GeV。在大多数计算中，$(a\mu) \approx 1$（甚至稍大），上界并未被严格满足。另一方面，该界限也依赖于所涉及作用量的标度性质。
由于 (11.63) 依赖规范，人们必须在固定规范下工作，并且必须在同一规范下比较依赖规范的格点矩阵元与连续统结果。朗道规范固定是一个合适的选择，但必须记住，Gribov 副本的不确定性可能会破坏比较。在格点计算中，人们发现这种效应的信号即使存在也微乎其微 [16--18]。
我们按照 [19] 对 [14] 的修改来总结这一方法。注意，在 (11.63) 中比较的是色空间和 Dirac 空间中的矩阵。取迹后，重整化条件为
\begin{equation}
\frac{Z_\Gamma}{12}\,
\tr\Bigl[\braket{p}{O_\Gamma}{p}\,\bigl(\braket{p}{O_\Gamma}{p}_0\bigr)^{-1}\Bigr]
\Big|_{p^2 = \mu^2} = 1 .
\end{equation}
矩阵元
\begin{equation}
\braket{p}{O_\Gamma}{p} = \frac{1}{Z_q}\,\Lambda_\Gamma(p)
\end{equation}
正比于截肢（amputated）Green 函数
\begin{equation}
\Lambda_\Gamma(p) = S^{-1}(p)\, G_\Gamma(p)\, S^{-1}(p) ,
\end{equation}
而 $Z_q$ 是将在下面讨论的夸克场重整化常数。Green 函数 $G_\Gamma(p)$ 作为期望值确定：
\begin{equation}
G_\Gamma(p)_{\alpha\beta}^{ab} =
\frac{1}{V}\sum_{x,y,z} e^{-ip(x-y)}
\ev{u_\alpha^a(x)\, O_\Gamma(z)\, \bar{d}_\beta^b(y)} .
\end{equation}
指标 $\alpha, \beta$ 和 $a, b$ 分别遍历 Dirac 指标和色指标，$V$ 表示格点体积。夸克传播子为
\begin{equation}
S_{\alpha\beta}^{ab}(x, y) = \ev{u_\alpha^a(x)\, \bar{u}_\beta^b(y)}
= \ev{d_\alpha^a(x)\, \bar{d}_\beta^b(y)}
\end{equation}
（假设 $u$ 和 $d$ 质量相等，并对期望值使用朗道规范）。
因此我们必须计算 $G_\Gamma(p)$ 和 $S(p)$。这按如下方式进行。对于在单个规范组态 $n$ 上求值的夸克传播子 $S^{(n)}$，我们定义
\begin{equation}
S^{(n)}(x|p) = \sum_y e^{ipy}\, S^{(n)}(x, y) .
\end{equation}
考虑到传播子的 $\gamma_5$-厄米性，对于如 (11.62) 中定义的夸克双线性算符 $O_\Gamma$，我们可以用 (11.70) 中的量改写 $G_\Gamma(p)$：
\begin{equation}
G_\Gamma(p) =
\frac{1}{V}\sum_{x,y,z} e^{-ip(x-y)}
\ev{u(x)\,\bar{u}(z)\,\Gamma\, d(z)\,\bar{d}(y)} ,
\end{equation}
\begin{equation}
G_\Gamma(p) \approx \frac{1}{VN}\sum_{n=1}^{N}\sum_z
\gamma_5\, S^{(n)}(z|p)^\dagger\, \gamma_5\, \Gamma\, S^{(n)}(z|p) ,
\end{equation}
其中我们通过对 $N$ 个规范组态求平均来近似期望值。类似地，我们得到动量空间中的夸克传播子
\begin{equation}
S(p) \approx \frac{1}{VN}\sum_{n=1}^{N}\sum_x e^{-ipx}\, S^{(n)}(x|p) .
\end{equation}
$S^{(n)}(y|p)$ 通过求解格点 Dirac 方程（$D$ 表示 Dirac 算符）来计算
\begin{equation}
\sum_y D(z, y)\, S^{(n)}(y|p) = e^{ipz}
\end{equation}
其中使用动量源（参见 [19]）。这有一个缺点，即必须为若干个动量源确定夸克传播子，而在原始方法 [14] 中，人们使用点源（即只考虑 $z = 0$ 而不是对所有 $z$ 求和）并将夸克汇点（sink）投影到所需的动量值。然而，使用动量源有一个很大的优点，即信号显著更好。
夸克场重整化常数通过将夸克传播子与自由（格点）传播子比较而得到。使用所谓 RI$'$ 方案，我们取
\begin{equation}
Z_q' =
\frac{1}{12}\,
\tr\left[S^{-1}(p)\,
\frac{R(p) - i\gamma_\nu v_\nu(p)}
{R(p)^2 + \sum_\nu v_\nu(p)^2}\right]
\Bigg|_{p^2 = \mu^2} .
\end{equation}
RI$'$ 方案与 RI 方案的区别仅在于夸克场重整化常数的定义。在 (11.75) 中，$R$ 和 $v_\nu$ 是自由 Dirac 算符中出现的标量和矢量项，后者在动量空间中为 $D(p) = i\gamma_\nu v_\nu(p) + R(p)$。利用傅里叶变换，可以从 $D$ 的定义计算出 $R$ 和 $v_\nu$。它们被归一化，使得
\begin{equation}
v_\nu(p) = i p_\nu + \order{ap^2} , \qquad
R(p) = \order{ap^2} .
\end{equation}
朗道规范固定 $\partial_\nu A_\nu = 0$ 按照 [17, 20] 中讨论的方式实现，即通过随机过松弛 [21] 迭代地最小化链变量的某个泛函（参见第 3.2.2 节）。众所周知，这种类型的规范固定仍然允许对应于进一步局部极小的 Gribov 副本。这会产生一种不确定性，必须对其进行检查，例如通过研究随机规范变换的效果。
容易验证，在树级有 $\braket{p}{O_\Gamma}{p}_0 = \Gamma$。综合起来，我们得到 RI$'$ 方案中 $Z_\Gamma$ 的最终公式：
\begin{equation}
Z_\Gamma^{\mathrm{RI}'} = \frac{12\,Z_q'}{\tr\bigl[\Lambda_\Gamma(p)\,\Gamma^{-1}\bigr]}
\Bigg|_{p^2 = \mu^2} .
\end{equation}
对于 $\Gamma = \gamma_\nu, \gamma_\nu\gamma_5$，隐含在求迹下对指标 $\nu$ 取平均。
不同重整化方案之间的联系利用连续统微扰论建立。在 $\overline{\mathrm{MS}}$ 和 RI$'$ 两个方案中，重整化因子相联系，我们得到格点数 $Z_\Gamma^{\mathrm{RI}'}$ 的最终结果
\begin{equation}
Z_\Gamma^{\overline{\mathrm{MS}}}(\mu^2) = R_\Gamma(\mu^2)\, Z_\Gamma^{\mathrm{RI}'}(\mu^2) ,
\end{equation}
其中比值 $R_\Gamma(\mu^2)$ 可以在微扰论中计算。$Z_\Gamma^{\mathrm{RI}'}(\mu^2)$ 和 $Z_\Gamma^{\overline{\mathrm{MS}}}(\mu^2)$ 在连续极限中都可能发散，但比值保持有限。
RI、RI$'$ 与 $\overline{\mathrm{MS}}$ 方案之间的转换因子已在朗道规范下、3 圈阶在 [15, 22, 23] 中确定。对重整化群微分方程积分，还可以定义与标度无关的量 $Z^{\mathrm{RGI}}$（RGI = 重整化群不变量）[19, 24]。
这里我们只讨论了重整化形如 (11.62) 的夸克双线性的简单情形。含有导数的算符，正如它们被用于第~\ref{sec:matrixelements} 节中讨论的矩阵元那样，可以用类似的方法处理。
在 RI 方案中，为了与质量无关的 $\overline{\mathrm{MS}}$ 方案比较，必须外推到手征极限。另一种方法是 Schrödinger 泛函方法（参见第 9.1 节），其中可以直接在手征点处工作 [13, 25]。

\section{强子衰变与散射}
\subsection{阈值区域}
在有限体积中，可观测量表现出有限体积效应。让我们首先讨论单个强子的传播。我们假设时间方向足够大，因此对有限体积效应不重要。我们还假设体积的空间方向采用周期性边界条件。那么单粒子传播子受有限体积影响有两种途径：一方面，在欧氏时间上相隔一定距离的两点之间的传播子实际上是源与所有空间镜像之间传播子的和。另一方面，质量重整化的相互作用（如圈）也会有路径沿一个周期性空间方向绕行一周再回到顶点的情形，如图 11.1 所示。这可以被解释为有限盒子对传播粒子的虚极化云的挤压效应。主导效应来自系统中最轻的相互作用粒子，即 QCD 中的$\pi$介子。控制参数是盒子大小 $L$ 与轻粒子关联长度之比，或者等价地，乘积 $M_\pi L$。

\begin{figure}[htbp]
  \centering
  \includegraphics[width=0.8\textwidth]{fig111.png}
  \caption*{图 11.1. \emph{$\pi$介子的传播。左：自由的（上方）和带相互作用圈的（下方）。右：在有限、周期的空间盒子中额外的项，由于周期性边界条件，$\pi$介子还可以绕环面传播}。}
  \label{fig:111}
\end{figure}

这已在 [26] 中针对强子的连续统量子场论进行了讨论。在 QCD 中，夸克和胶子是基本对象。然而，由于禁闭，夸克圈被压低，预期主导效应来自$\pi$介子。因此这一推导应当适用于 QCD。
如图 11.1 所示，单个强子在发射和吸收$\pi$介子时不改变其动量。因此，这一相互作用与前向散射振幅相联系。如果在某个近似中，这只是一个常数 $\alpha$，那么质量改变正比于（发散的）圈贡献
\begin{equation}
m^2 = m_0^2 + \alpha \int \frac{d^4p}{(2\pi)^4}\, G(p) = m_0^2 + \alpha\, G(0) ,
\end{equation}
其中 $G(p)$ 是动量空间中的传播子。积分给出原点处的实空间传播子 $G(0)$，正如 (11.79) 第二步所明确的。发散被吸收进裸质量的定义中（重整化）。
现在我们考虑大小为 $L^3$ 的有限、周期空间盒子，其时间延伸足够大，可以忽略由有限性带来的修正。在这样的有限体积中，我们必须把实空间传播子替换为求和
\begin{equation}
G(x) \Rightarrow \sum_{s \in \mathbb{Z}^3} G(x + L s) ,
\end{equation}
其中 $s \in \mathbb{Z}^3$ 给出原始格点与其镜像之间的平移。
由相互作用项，即由来自镜像传播的贡献引起的有限与无限盒子之间的质量差为
\begin{equation}
m^2(L) - m^2 = \alpha \sum_{s \neq 0} G(L s) .
\end{equation}
为了估计这些贡献，我们现在需要自由玻色子传播子的大距离行为。有多种方式可以确定傅里叶变换
\begin{equation}
G(x) = \int \frac{d^4p}{(2\pi)^4}\, \frac{e^{ipx}}{p^2 + m^2}
\sim \frac{m^2 \sqrt{8\pi}\, e^{-m|x|}}{(4\pi)^2 (m|x|)^{3/2}} .
\end{equation}
因此我们可以预期主导的体积依赖 $\propto e^{-mL}/L^{3/2}$。系数可以例如从 ChPT [27--29] 确定。
Lüscher [26] 对周期、有限格点上圈贡献的微扰论所有阶求和，得到表达式
\begin{equation}
m(L) - m = -\frac{1}{2}\,\frac{3}{(4\pi)^2 mL}
\int_{-\infty}^{\infty} dy\, e^{-\sqrt{y^2 + m^2}\, L}\, F(iy) + O(e^{-mL}) ,
\end{equation}
其中第二项被压低，因为可以证明 $m > \tfrac{3}{2}m$。$F(iy)$ 表示 $X\pi \to X\pi$ 前向散射振幅 $F(\nu = s - u) = A(s, t = 0, u)$ 解析延拓到虚宗量的形式，其中 $s$、$t$ 和 $u$ 是相对论不变量的 Mandelstam 变量。
我们这里不重复推导，只是使结果看起来合理。为此，我们沿时间轴选取 $x$ 来积分 (11.82)，并写出
\begin{equation}
\begin{aligned}
G(x) &= \int_{-\infty}^{\infty} \frac{dE}{2\pi}\, e^{iEx}
\int_{-\infty}^{\infty} \frac{d^3p}{(2\pi)^3}\, \frac{1}{E^2 + \mathbf{p}^2 + m^2} \\
&= \int_{-\infty}^{\infty} \frac{dE}{2\pi}\, e^{iEx}
\int_0^{\infty} \frac{dy}{2\pi}\, \frac{y^2}{E^2 + y^2 + m^2}
\end{aligned}
\end{equation}
（这里 $y = |\mathbf{p}|$）。对 $E$ 的积分可以用柯西定理求解，即在无穷远处用半圆闭合围道，并取上半平面极点 $E = i\sqrt{y^2 + m^2}$ 处的留数，
\begin{equation}
\begin{aligned}
G(x) &= \frac{1}{4\pi^2} \int_0^{\infty} dy\,
\frac{y^2}{\sqrt{y^2 + m^2}}\, e^{-\sqrt{y^2 + m^2}\, x} \\
&= -\frac{1}{4\pi^2}\,\frac{1}{x} \int_0^{\infty} dp\, p\,
\frac{\partial}{\partial p}\, e^{-\sqrt{p^2 + m^2}\, x}
= \frac{1}{8\pi^2 x} \int_{-\infty}^{\infty} dp\, e^{-\sqrt{p^2 + m^2}\, x} .
\end{aligned}
\end{equation}
在第二步中我们做了分部积分，然后加倍积分区间，以得到与 (11.83) 类似的形式。我们现在只把求和 (11.81) 中来自六个最近镜像的领头贡献相加，得到
\begin{equation}
m^2(L) - m^2 = \frac{6\alpha}{8\pi^2 L}
\int_{-\infty}^{\infty} dp\, e^{-\sqrt{p^2 + m^2}\, L} ,
\end{equation}
由此
\begin{equation}
m(L) - m \approx \frac{m^2(L) - m^2}{2m}
= \frac{3\,(-2\alpha)}{2m}\,
\frac{1}{(4\pi)^2 L}
\int_{-\infty}^{\infty} dp\, e^{-\sqrt{p^2 + m^2}\, L} .
\end{equation}
这正是 (11.82)，其中散射振幅被替换为常数耦合 $F \to -2\alpha$。这样我们就重现了前面讨论过的简单领头行为。
在该系列的第二篇论文 [30] 中，Lüscher 随后将有限周期盒子中两个稳定粒子（静止）的能谱与其弹性散射振幅联系起来。单粒子的有限体积效应是指数小的，而现在两个粒子之间的相互作用只导致多项式压低 $O(1/L^3)$。其物理图景是，局域的粒子需要彼此靠近才能相互作用，而找到一个粒子靠近另一个粒子的概率与空间体积成反比。
在阈值附近，散射振幅由 $s$ 波散射长度 $a_0 = \lim_{k\to 0} \delta_0(k)/k$ 主导（其中 $k$ 是二粒子动量，$\delta_0$ 是 $s$ 波散射相移）。对于二粒子系统，基态能量按 $1/L$ 的展开为
\begin{equation}
W = 2M_\pi - \frac{4\pi a_0}{M_\pi L^3}
\left[1 + c_1\frac{a_0}{L} + c_2\frac{a_0^2}{L^2} + O\left(\frac{1}{L^3}\right)\right] .
\end{equation}
系数 $c_1 = -2.837297$ 和 $c_2 = 6.375183$ 依赖于格点形状，并与动量格点的广义 zeta 函数相关。针对非零动量的类似关系还涉及其他相移 $\delta_l(k)$。
关系 (11.88) 可以推广到质量不等的粒子以及非零总动量的情形。领头项始终与体积的倒数成正比 [30]。

\subsection{阈值区域之外}
渐近地只能观察到稳定态。因此，在全 QCD 模拟中，共振态通常必须通过它们对中间态的影响来识别。利用如 (11.88) 的方程，Lüscher [26, 30--32] 推导出有限格点上二粒子本征态能量与散射相移之间的关系，该关系在非弹性阈值之下有效。
让我们在一个简单的情形中讨论这一思想。考虑一个空间延伸为有限 $L$ 但时间延伸无限的一加一维系统中两个玻色子的散射。这一情形最方便地用散射伙伴相对运动的波函数来描述。基本假设是相互作用程是有限的且小于 $L$。因此，相互作用程之外的波函数是平面波。相互作用的效果由一个动量依赖的相移 $\delta(k)$ 来考虑，该相移在相互作用区域中获得，如图 11.2 所示。

\begin{figure}[htbp]
  \centering
  \includegraphics[width=0.8\textwidth]{fig112.png}
  \caption*{图 11.2. \emph{该图说明波函数的行为：在相互作用区域之外，它是未受扰动的平面波，在相互作用区域内获得额外的相移（用箭头表示）}。}
  \label{fig:112}
\end{figure}

施加周期性边界条件，动量 $k$ 的平面波在 $x = L$ 处的匹配意味着
\begin{equation}
e^{ikL + 2i\delta(k)} = e^{ik\cdot 0} = 1 .
\end{equation}
这给出相对动量 $k_n$ 的量子化条件：
\begin{equation}
k_n L + 2\delta(k_n) = 2n\pi ,
\end{equation}
其中 $n$ 是整数。当 $\delta$ 为零时，回到自由情形，此时动量为 $k_n = 2\pi n/L$。对于相互作用情形，可以从 (11.90) 推出非平凡的动量。如果函数形式 $\delta(k_n)$ 已知，就可以利用这一关系找出该有限体积中动量的量子化值。另一方面，给定来自某种测量的动量谱，(11.90) 允许确定每个 $k_n$ 的相移 $\delta(k_n)$。
动量可以从模拟中可获得的二粒子态的能量值得到。对于给定的 $L$，人们计算分立能级 $W_0, W_1, W_2, \ldots$，并利用色散关系从中得到 $k_n$ 的值
\begin{equation}
W_n = 2\sqrt{m^2 + k_n^2} .
\end{equation}
技术问题在于精确确定单粒子质量和能级 $W_n$。
为了确定能谱，必须使用诸如变分法（在第 6.3.3 节中讨论）之类的技术，考虑足够多数量的、具有正确量子数且能够表示散射态空间 [33, 34]（包括耦合的单粒子道）的插入算符的关联函数。通常，若干最低能量的本征模可以以足够的可靠性确定。改变系统的空间大小 $L$ 使得人们能够覆盖不同的动量值。在 [34] 中研究了一个简单的二维系统，它在格点上耦合一个较重和两个较轻的玻色子，其质量和耦合参数允许发生类似 $\rho \to \pi\pi$ 系统的衰变。图 11.3 展示了预期的能级回避交叉现象，它导致共振的相移。

\begin{figure}[htbp]
  \centering
  \includegraphics[width=0.9\textwidth]{fig113.png}
  \caption*{图 11.3. \emph{在一个将一个重粒子耦合到两个轻粒子的简单玩具模型中，演示散射相移与有限体积能谱之间的关系（图形来自 [34]）。我们显示了在没有 (a) 和有 (b) 将单个重粒子耦合到两个轻粒子的相互作用情况下观察到的能级。(a) 和 (b) 中的水平虚线表示 2 粒子和 4 粒子阈值，(a) 中的水平实线表示单个（重）粒子质量。图 (c) 显示 (b) 情形下产生的共振相移；竖直虚线表示 4 粒子阈值。对于二维模型，相移从 $-\pi/2$ 开始（无相互作用）。在所有图中，符号表示数值数据，实曲线表示自由情形 (a) 或共振相互作用情形 (b)、(c) 的模型}。}
  \label{fig:113}
\end{figure}

必须仔细考虑该方法的局限性：相互作用区域和单粒子关联长度应当小于空间体积，特别是 $mL \gg 1$。该关系只适用于第一个非弹性阈值之下。由虚粒子绕环面传播引起的极化效应应当处于可控范围。对于大的 $k$ 值，会出现格点人为效应。
在物理的 4 维情形中，相移、格点大小与动量之间的关系变得稍微复杂：
\begin{equation}
\delta(k) = \phi(q) \bmod \pi , \qquad
q = \frac{kL}{2\pi} , \qquad
\tan(-\phi(q)) = \frac{q\,\pi^{3/2}}{\mathcal{Z}_{00}(1; q^2)} , \qquad
\phi(0) = 0 .
\end{equation}
这里 $\mathcal{Z}_{00}$ 表示黎曼 zeta 函数。一个 4 维玻色子模型已在 [35, 36] 中研究，一个淬火 QCD 模拟在 [37] 中研究。
$\rho$ 介子是 $\pi\pi$ $p$ 波道中的共振态，因此衰变中的$\pi$介子具有非零动量。在 [38] 中（另见 [39]），该关系被推广到非静止系情形。其优点是，在同一个组态集合上可以同时研究零动量和非零动量两种情形，从而得到更多的数据点。可以利用更小的体积和更小的体积范围。一个关于 $\rho$ 衰变的动力学费米子研究使用了这一方法 [40, 41]。
在 [42] 中，通过在 $\rho$ 质量位于两$\pi$介子质量之下的情形中分析能级来研究衰变。从 $N_T$ 依赖性确定了跃迁振幅 $\rho \to \pi\pi$，并由此确定有效耦合。

\section{矩阵元}
\label{sec:matrixelements}
在探索强子的内部结构时，人们想知道电荷、物质和自旋的空间分布。然而，人们还需要知道动量分布。这两个量本质上都是非微扰的，因此对格点 QCD 是一个挑战。
空间分布可以通过弹性散射过程来探测，由此引出形状因子和电荷分布的概念。动量分布可以在诸如深度非弹性散射或 Drell–Yan 过程之类的硬（高能）碰撞中通过实验来探索。

\begin{figure}[htbp]
  \centering
  \includegraphics[width=0.8\textwidth]{fig114.png}
  \caption*{图 11.4. \emph{介子态（圆圈）之间算符（方块）的矩阵元}。}
  \label{fig:114}
\end{figure}

在部分子模型 [43] 中，夸克、反夸克和胶子都是部分子，并在具有很高动量的强子中被考虑。点状探针确定部分子概率分布，即在一个具有强子纵向动量某一份额 $x$ 和动量转移 $Q^2$ 的无穷动量参考系中找到部分子的概率分布。
强子 $h$ 中的夸克分布函数通过连续统光锥记法中的一个矩阵元来定义
\begin{equation}
q(x) = \frac{1}{2p^+}\int \frac{d\lambda}{2\pi}\, e^{i\lambda x}\,
\ev{h(p)}{\bar{\Psi}(0)\gamma^+\Psi(\lambda n)}{h(p)} ,
\end{equation}
其中 $p^+ = (p^0 + p^3)/\sqrt{2}$，$n^\mu = (1, 0, 0, -1)$，且 $n \cdot p = 1$。夸克场 $\Psi$ 用规范输运线从其所处位置连接到无穷远处，以使该表达式具有规范不变性。
广义部分子分布（GPD）是这一概念的推广，引入诸如自旋和横向动量之类的进一步变量，并将空间值与动量值关联起来。通过这一推广，人们希望也能描述较低能量下的散射过程 [44, 45]。
威尔逊的算符乘积展开（OPE）[46--49] 将光锥矩阵元 (11.93) 或 GPD 与局域算符 $\ev{h}{O_i}{h'}$ 的矩阵元联系起来。还需要重整化因子 $Z_{O_i}(a\mu)$ 将格点结果转换到例如 $\overline{\mathrm{MS}}$ 方案。局域算符可以是简单的夸克双线性，如 $\bar{u}\gamma_\nu d$（用于形状因子），它们可以包含导数，如 $\bar{u}\gamma_\nu D_\mu d$（用于 GPD），或者包含更多夸克，如 $\bar{s}\gamma_\mu u\,\bar{u}\gamma_\mu d$（用于弱相互作用过程）。图 11.4 给出了需要确定的矩阵元的例子。

\subsection{$\pi$介子形状因子}
介子形状因子是最简单的矩阵元之一。作为原型例子，我们讨论描述$\pi$介子与光子耦合的$\pi$介子电磁形状因子 $F^{(\pi)}$。它由矩阵元定义
\begin{equation}
\ev{\pi^+(p_f)}{V_\mu}{\pi^+(p_i)}_{\mathrm{cont}} = (p_f + p_i)_\mu\, F^{(\pi)}(Q^2) ,
\end{equation}
其中 $p_i$ 和 $p_f$ 分别是初态和末态$\pi$介子的 4 动量，$Q^2 = (p_f - p_i)^2 \equiv -t$ 是类空的不变动量转移平方，而
\begin{equation}
V_\mu = \frac{2}{3}\,\bar{u}\gamma_\mu u - \frac{1}{3}\,\bar{d}\gamma_\mu d
\end{equation}
是同位旋矢量流，在零动量转移处归一化为单位电荷（即在壳光子耦合）：$F^{(\pi)}(0) = 1$。它的类空值由实验确定。
与连续统归一化的联系给出另一个因子，
\begin{equation}
\ev{\pi^+(p_f)}{V_\mu}{\pi^+(p_i)}_{\mathrm{latt}}
= \frac{1}{\sqrt{2E_iE_f}}\,
\ev{\pi^+(p_f)}{V_\mu}{\pi^+(p_i)}_{\mathrm{cont}} .
\end{equation}
$\pi$介子形状因子是 $t$ 的解析函数。它的类空值可以由其解析性区域边界上的值确定，即沿正实 $t$ 轴的割线，即类时区域。虽然割线从 $t = 4m_e^2$ 开始，但显著的贡献只来自强子共振区域，从 $4M_\pi^2$ 处的两$\pi$介子阈值开始。由于割线前段（直到非弹性道变得重要为止）的幺正性，相移本质上就是弹性两$\pi$介子道的 $p$ 波相移。该区域由 $\rho$ 介子共振主导。这一特征已被各种色散关系表示所利用 [50--53]。
尽管四$\pi$介子道还有进一步的非弹性贡献，但这些贡献在 $\pi\omega$ 道开启之前一直很小。还存在通过高阶电磁相互作用的同位旋标量 $\omega$ 耦合的小贡献。所有这些贡献在近类空区域都只显示很小的影响。矢量介子的主导地位促使人们对这一区域采用简单的极点参数化。
平均电荷半径平方 $\ev{r^2}_V$ 通过下式定义
\begin{equation}
F^{(\pi)}(Q^2) = 1 - \frac{1}{6}\,\ev{r^2}_V\, Q^2 + O(Q^4)
\quad\Rightarrow\quad
\ev{r^2}_V \equiv 6\,\frac{dF^{(\pi)}(t)}{dt}\Big|_{t=0} .
\end{equation}
其值的当前 PDG 平均值为 $0.45(1)~\mathrm{fm}^2$ [2]。
利用现有的格点工具，人们在类空区域研究$\pi$介子形状因子。为此，我们需要计算离前向矩阵元（最好在几个转移动量处），即如下形式的期望值
\begin{equation}
\sum_{y;\,y_0=\tau}\;\sum_{x;\,x_0=t}
\tr\Bigl[e^{iq\cdot y}\, S(0, y)\, O(y)\,
e^{-ip\cdot x}\, S(y, x)\,\gamma_5\, S(x, 0)\,\gamma_5\Bigr] ,
\end{equation}
其中 $S(y, x)$ 是从 $x$ 到 $y$ 的夸克传播子，$O(y)$ 表示在 $y$ 处插入的算符。我们使用记号 $p \equiv p_f$，$r \equiv p_i$，并把动量转移记为

\begin{figure}[htbp]
  \centering
  \includegraphics[width=0.8\textwidth]{fig115.png}
  \caption*{图 11.5. \emph{矩阵元 (11.98) 的示意图}。}
  \label{fig:115}
\end{figure}

$q = p - r$（参见图 11.5）。理想情况下，源与算符插入之间以及算符插入与汇点之间的时间距离足够长，使得$\pi$介子道中的更高激发可以被忽略，$\pi$介子基态主导插入区域。对源和汇点算符进行涂抹（smearing）也有助于获得更好的基态信号。
在计算矩阵元时，必须计算如图 11.4 和图 11.5 中那样的夸克传播子组合的期望值。每一次从给定源计算夸克传播子都需要对 Dirac 算符求逆，从而需要大量的计算时间。因此，人们希望减少传播子计算的次数。这通过使用所谓的顺序源方法 [54, 55] 来实现。我们首先如第 6 章所述，计算从原点处的源到格点所有其他点的夸克传播子。这同时提供了从源原点传播到 $x$ 的传播子，并且由于 $\gamma_5$-厄米性，还提供了从 $y$ 传播到原点的传播子。我们所需要的只是一个从 $x$ 到 $y$ 的传播子。矩阵元于是写为
\begin{equation}
\sum_{y;\,y_0=\tau}
\tr\Bigl[e^{iq\cdot y}\, S(0, y)\, O(y)\, \Sigma(y, 0)\, \gamma_5\Bigr] ,
\end{equation}
其中顺序传播子定义为
\begin{equation}
\Sigma(y, 0) = \sum_{x;\,x_0=t} e^{-ip\cdot x}\, S(y, x)\,\gamma_5\, S(x, 0) .
\end{equation}
对于末态动量 $p$ 的每一种选择，它可以通过对 Dirac 算符 $D$ 进行额外的一次求逆来容易地计算：
\begin{equation}
\sum_y D(z, y)\, \Sigma(y, 0) = e^{-ip\cdot z}\, \gamma_5\, S(z, 0) , \qquad z_0 = t .
\end{equation}
改变汇点的性质需要计算新的顺序传播子，因此模拟多个末态动量、不同的场插入算符或汇点的不同涂抹会很快变得相当昂贵。出于这个原因，人们通常只在一个末态动量值处工作。这种方法不仅可以确定所提到的矩阵元，而且不费额外力气，还可以考虑其他算符插入，例如用于标量形状因子。
人们通过计算三点与两点关联函数之比来提取物理矩阵元。为了说明这一点，让我们考察入射和出射赝标量插入场算符 $P$ 以及一个在插入时刻 $\tau$ 处的算符 $O$ 的三点关联函数。在 $\tau$ 处我们插入一组完备态。当分离距离足够大时，这可以被该道中的基态——$\pi$介子——所饱和，我们得到
\begin{equation}
\begin{aligned}
\ev{0}{P(t; p)\, O(\tau)\, P(r)}{0}
&= \ev{0}{P(t; p)}{\pi(p)}\,
\frac{e^{-E_p(t-\tau)}}{2E_p}\\
&\qquad \ev{\pi(p)}{O(\tau)}{\pi(r)}\,
\frac{e^{-E_r\tau}}{2E_r}\,
\ev{\pi(r)}{P(0; r)}{0} ,
\end{aligned}
\end{equation}
而对于赝标量传播子，
\begin{equation}
\ev{0}{P(t; p)\, P(p)}{0}
= \ev{0}{P(t; p)}{\pi(p)}\,
\frac{e^{-E_p t}}{2E_p}\,
\ev{\pi(p)}{P(0; p)}{0} .
\end{equation}
这里我们忽略时间方向的格点周期性，否则将会有来自反向传播的进一步项。这一点已例如在 [56] 中讨论过。
现在构造这些函数的合适组合，使得未知因子和指数项相互抵消：
\begin{equation}
\begin{aligned}
R(t, \tau; p, q \equiv p - r)
&= \sqrt{2E_iE_f}\,
\frac{\ev{P(t; p)\, O(\tau; q)\, P(0; r)}}{\ev{P(t; p)\, P(0; p)}} \\
&\times
\sqrt{
  \frac{\ev{P(t; p)\, P(0; p)}}{\ev{P(t; r)\, P(0; r)}}
  \frac{\ev{P(\tau; p)\, P(0; p)}}{\ev{P(\tau; r)\, P(0; r)}}
  \frac{\ev{P(t - \tau; r)\, P(0; r)}}{\ev{P(t - \tau; p)\, P(0; p)}}} \\
&= \ev{\pi(p)}{O(\tau; q)}{\pi(r)} .
\end{aligned}
\end{equation}
可以将汇点固定在时刻 $t$，并改变算符 $O$ 所在的时间切片 $\tau$（扫描一系列时间切片）。因此，$R(t, \tau; p, q)$ 在 $\tau$ 上呈现两个平台：$0 \ll \tau \ll t$ 和 $t \ll \tau \ll N_T$。当汇点放在 $t = N_T/2$ 时，该比值（在我们的情形中）关于 $\tau - N_T/2$ 是反对称的。形状因子的值可以从平台值读出。由于源和汇点算符涂抹引起的进一步归一化因子在比值中也相互抵消。
$Q^2 = 0$ 时得到的平台值不一定是 1，即我们例子中的电荷。这是因为 (11.95) 中朴素的格点流 $V_\mu$ 并不守恒。人们可以用守恒流来替换它，然而对于手征作用量，这在数值上要求很高，正如第~\ref{sec:naivecurrents} 节所讨论的。对于其他作用量，即使守恒流也需要 $O(a)$ 修正项才能得到改进的流（参见第 9.1 节）。采用简单的点状流时，必须确定将其与重整化流 $V_\mu^{(r)}$ 相联系的重整化因子 $Z_V$。
源、插入和汇点位于三个时间切片上。源算符通常局域在一个格点或某个格点附近。在傅里叶分解中，它因此对所有可能的 3 动量都有贡献。通过将汇点投影到 3 动量 $p$、插入投影到 $q$，源的相关分量为 $r = p - q$。由于平移不变性，所有其他贡献在对所有场组态求和时相互抵消。在具有大量但有限的规范组态的现实情形中，这种抵消只是近似的。
对于$\pi$介子电磁形状因子，合适的 3 动量选择是 $|p_f| = |p_i|$，使得 $E_i = E_f$。在这种情况下 $Q^2 = |\mathbf{q}|^2$，结合 (11.94) 和 (11.96)，运动学因子因下式而简化
\begin{equation}
\frac{E_i + E_f}{\sqrt{2E_iE_f}} = 1 .
\end{equation}
这消除了另一个可能的统计误差来源。一种最小化计算量的可行方法是将汇点的动量固定为 $p = (1, 1, 0)\,\Delta p$，其中 $\Delta p = 2\pi/(aN)$ 是空间动量间隔。这样，对每个组态只需额外计算一个（顺序）传播子。然后对满足 $|r| = |p - q| = \sqrt{2} = |p|$ 的 $q$ 值求三点函数。这些是 12 个组合 $(0, 0, 0)$、$(2, 2, 0)$、$(2, 0, 0)$、$(0, 2, 0)$、$(0, 1, \pm 1)$、$(1, 0, \pm 1)$、$(2, 1, \pm 1)$、$(1, 2, \pm 1)$。这给出 $Q^2 = 2n(\Delta p)^2$（$n = 0, 1, 2, 3, 4$）的结果。
两点和三点函数的动量投影会产生显著的统计噪声，特别是在对称点 $t = N_T/2$ 附近。因此，人们把汇点的位置选在离源稍近一些的时间切片上。矩阵元随后通过组合汇点两侧的平台值 $R_{\mathrm{lhs}} \pm R_{\mathrm{rhs}}$ 获得，其中相对符号取决于插入算符的宇称性质。对于矢量形状因子，符号为负。
$\pi$介子形状因子最早的格点研究是 [57, 58]。类空形状因子通常被类时 $\rho$ 极点很好地近似，而在非物理的大夸克质量下的淬火计算 [59, 60] 中，它大于实验值。对于具有动力学夸克并外推到物理夸克质量的模拟，与物理 $\rho$ 质量和所测$\pi$介子形状因子的一致性得到改善 [56, 61, 62]。$\pi$介子处于手征对称性破缺的核心，因此 ChPT 提供了若干约束 [53, 63--65]。

\subsection{弱相互作用矩阵元}
弱相互作用几乎以点状方式与夸克耦合，但夸克生活在强子的 QCD 主导环境中。非微扰效应决定了强子结构，因此对于分析弱相互作用“看待”强子的方式很重要。
在标准模型（SM）中，电弱相互作用由光子和重矢量玻色子 $W^\pm$ 和 $Z^0$ 介导。这些玻色子与轻子的双线性组合和夸克的双线性组合耦合，从而允许夸克味改变过程。六种夸克味排列在三个代（family）中：$(u, d)$、$(c, s)$ 和 $(t, b)$。弱相互作用“看待”夸克味的方式与强相互作用略有不同。相应的味本征态记为 $d'$、$s'$ 和 $b'$，它们通过 Cabibbo–Kobayashi–Maskawa（CKM）矩阵 [66, 67] 与质量本征态 $d$、$s$ 和 $b$ 相联系
\begin{equation}
\begin{pmatrix} d' \\ s' \\ b' \end{pmatrix}
= V_{\mathrm{CKM}} \begin{pmatrix} d \\ s \\ b \end{pmatrix}
= \begin{pmatrix}
V_{ud} & V_{us} & V_{ub} \\
V_{cd} & V_{cs} & V_{cb} \\
V_{td} & V_{ts} & V_{tb}
\end{pmatrix}
\begin{pmatrix} d \\ s \\ b \end{pmatrix} .
\end{equation}
在不失一般性的情况下，其余三种夸克可以保持不变。CKM 矩阵是复的（允许 CP 破坏）且幺正的（在树级强制不存在味改变中性流跃迁）。关于这些方面的综述，参见 [68, 69]。
在若干进一步的、由唯象学动机驱动的近似下，该矩阵通常用四个所谓的广义 Wolfenstein 参数来表述：$\lambda = V_{us} = \sin\theta_c$、$A$、$\rho$ 和 $\eta$。这些参数可以与各种实验相联系。非零的 $\eta$ 是标准模型中 CP 破坏的原因。
CKM 矩阵的幺正性给出如下关系
\begin{equation}
V_{ud}V_{ub}^* + V_{cd}V_{cb}^* + V_{td}V_{tb}^* = 0 .
\end{equation}
这一关系可以在 $(\rho, \eta)$ 平面中表示为“幺正三角形”。三角形底边的端点是 $(0, 0)$ 和 $(1, 0)$，第三个角随后受到实验的约束，特别是 $K$ 和 $B$ 物理结果。格点计算有助于确定与 CKM 矩阵元计算相关的 QCD 输入 [70--72]。

\subsection{OPE 展开与有效弱哈密顿量}
\label{sec:weak}
强子的弱衰变涉及两个质量标度：重中间矢量玻色子 $W^\pm$ 和 $Z$，以及主导强子结构的 QCD 低能标度。人们通过引入一个有效弱哈密顿量将问题分成不同的部分
\begin{equation}
H_{\mathrm{eff}} = \frac{G_F}{\sqrt{2}} \sum_i V_{\mathrm{CKM}}\, C_i(\mu)\, Q_i(\mu) .
\end{equation}
这里 $G_F$ 是费米常数，$Q_i$ 表示由夸克场和轻子场构成的局域算符。CKM 因子和威尔逊系数 $C_i$ 决定算符的权重。这就是所谓的算符乘积展开（OPE）[46--49]。其中的一个局域算符正是来自 $\beta$ 衰变费米理论的著名的 $(V - A)$ 流乘积。这些项在强子态之间的期望值随后与实验可观测量相联系。
威尔逊系数（短距离标度）可以在微扰论中计算（参见 NLO 计算 [68]）。长距离（非微扰）标度被编码在重整化标度 $\mu$ 处求值的 $Q_i$ 中。通常 $\mu$ 的量级为 1--2~GeV 或更高；然而，它远小于弱相互作用标度 $M_W$，因此修正 $O(\mu^2/M_W^2)$ 可以忽略。在格点计算中，人们计算强子矩阵元 $\ev{h}{Q_i}{h'}$，即四夸克算符在强子态之间的矩阵元（参见图 11.4）。
当使用不满足 Ginsparg–Wilson 方程的格点 Dirac 算符时，手征对称性的缺失会导致显著的复杂性 [73]：
\begin{itemize}
\item 对于 $K - \bar{K}$ 混合和 $B_K$ 的计算，需要如 $O_{LL} = \bigl(\bar{s}\gamma_\mu(1 - \gamma_5)d\bigr)^2$ 的算符。如果手征对称性得不到保证，重整化时这些算符会与其他不同手征性的算符混合，如 $O_{PP} = \bigl(\bar{s}\gamma_5 d\bigr)^2$。矩阵元 $\ev{K}{O_{LL}}{\bar{K}}$ 在手征极限中消失（正比于夸克质量），而矩阵元 $\ev{K}{O_{PP}}{\bar{K}}$ 趋近于一个常数，因此需要高精度来分离出物理上有意义的部分。
\item 有效哈密顿量的算符会与更低维算符混合。一个例子是与 $K \to \pi\pi$ 相关的 $\Delta S = 1$ 哈密顿量。维度-6 算符 $\bar{s}\gamma_\mu(1 - \gamma_5)u\;\bar{u}\gamma_\mu(1 - \gamma_5)d$ 与如 $\bar{s}d$、$\bar{s}\gamma_5 d$ 和 $\bar{s}\gamma_\mu\gamma_\nu G_{\mu\nu} d$ 的算符混合（其中 $G$ 表示胶子场）。混合系数按 $1/a$ 的幂次发散，必须采用精确的非微扰方法来处理这些问题。
\end{itemize}

对于诸如威尔逊作用量这样破坏手征对称性的作用量，由于需要减除发散项，计算因此很繁琐。另一方面，使用 Ginsparg–Wilson 费米子会显著改善这一状况。对于 $\Delta I = \tfrac{1}{2}$ 问题，没有任何幂次发散的减除 [74]，而对于 $\varepsilon'/\varepsilon$ 的计算则需要减除（综述参见例如 [13]）。
非轻子弱衰变，其中末态完全由强子构成，如 $K \to \pi\pi$，是一项具有挑战性的任务。问题主要有两大类。紫外问题是从裸格点算符构造重整化算符的有限矩阵元；这在 [6, 75--77] 中讨论。红外问题与出射二粒子系统的再散射、从欧几里得到闵可夫斯基时空的延拓以及空间体积的效应有关。
Maiani–Testa 定理 [78] 指出，在非常大的空间体积中的欧几里得计算中，矩阵元被零动量单粒子态对主导。在无穷体积的欧几里得关联函数中，人们得到进入和离开两$\pi$介子态的矩阵元的平均值。然而，在处理有限体积时，这一状况会得到改善。Lellouch 和 Lüscher [79] 推导出有限体积中 $K \to \pi\pi$ 矩阵元与物理 K 介子衰变振幅之间的关系。在可达到的有限体积中，两$\pi$介子能谱远非连续的，静止的 K 介子不能衰变成两个$\pi$介子（除非某个二粒子能级接近其质量）。一个简单的表达式随后将有限体积中的跃迁振幅与无穷体积中的衰变率联系起来。该工作已被改编 [80]，以包含非弹性阈值以下的所有弹性态。这些方法还允许在 K 介子质量不等于两$\pi$介子能量时提取衰变振幅，即当插入的（弱哈密顿量）算符携带非零能量动量时。关于 K 介子物理格点结果的综述参见 [81]，更一般地，关于格点 QCD 能为标准模型做些什么，参见 [82]。

\section*{References}
\begin{enumerate}[label={[\arabic*]}]
  \item M. L\"uscher: in \emph{Probing the Standard Model of Particle Interactions}, Proceedings Les Houches 1997, edited by R. Gupta, A. Morel, E. de Rafael, and F. David, Vol. 1 (Elsevier, Amsterdam 1998) 268
  \item C. Amsler et al.\ [Particle Data Group]: Phys.\ Lett.\ B 667, 1 (2008) 270, 291
  \item M. Gell-Mann, R. J. Oakes, and B. Renner: Phys.\ Rev.\ 175, 2195 (1968) 274
  \item S. Weinberg: Physica A 96, 327 (1979) 274
  \item J. Gasser and H. Leutwyler: Nucl.\ Phys.\ B 250, 465 (1985) 274
  \item M. Bochicchio et al.: Nucl.\ Phys.\ B 262, 331 (1985) 274, 296
  \item P. Hasenfratz: Nucl.\ Phys.\ B 525, 401 (1998) 274
  \item P. Hasenfratz et al.: Nucl.\ Phys.\ B 643, 280 (2002) 275, 276, 277, 278
  \item Y. Kikukawa and A. Yamada: Nucl.\ Phys.\ B 547, 413 (1999) 276
  \item S. Capitani et al.: Phys.\ Lett.\ B 468, 150 (1999) 277
  \item L. Giusti, C. Hoelbling, and C. Rebbi: Phys.\ Rev.\ D 64, 114508 (2001) 277
  \item C. Gattringer, P. Huber, and C. B. Lang: Phys.\ Rev.\ D 72, 094510 (2005) 279
  \item S. Capitani: Phys.\ Rept.\ 382, 113 (2003) 280, 281, 284, 296
  \item G. Martinelli et al.: Nucl.\ Phys.\ B 445, 81 (1995) 281, 282, 283
  \item E. Franco and V. Lubicz: Nucl.\ Phys.\ B 531, 641 (1998) 281, 284
  \item M. L. Paciello, S. Petrarca, B. Taglienti, and A. Vladikas: Phys.\ Lett.\ B 341, 187 (1994) 282
  \item L. Giusti et al.: Int.\ J.\ Mod.\ Phys.\ A 16, 3487 (2001) 282, 283
  \item L. Giusti, S. Petrarca, B. Taglienti, and N. Tantalo: Phys.\ Lett.\ B 541, 350 (2002) 282
  \item M. G\"ockeler et al.: Nucl.\ Phys.\ B 544, 699 (1999) 282, 283, 284
  \item H. Suman and K. Schilling (unpublished) arXiv:hep-lat/9306018 (1993) 283
  \item J. E. Mandula and M. Ogilvie: Phys.\ Lett.\ B 48, 156 (1990) 283
  \item K. G. Chetyrkin and A. Retey: Nucl.\ Phys.\ B 583, 3 (2000) 284
  \item J. A. Gracey: Nucl.\ Phys.\ B 662, 247 (2003) 284
  \item V. Gimenez, L. Giusti, F. Rapuano, and M. Talevi: Nucl.\ Phys.\ B 531, 429 (1998) 284
  \item S. Capitani, M. L\"uscher, R. Sommer, and H. Wittig: Nucl.\ Phys.\ B 544, 669 (1999) 284
  \item M. L\"uscher: Commun.\ Math.\ Phys.\ 104, 177 (1986) 284, 286, 287
  \item P. Hasenfratz and H. Leutwyler: Nucl.\ Phys.\ B 343, 241 (1990) 286
  \item A. Hasenfratz et al.: Z.\ Physik C 46, 257 (1990) 286
  \item A. Hasenfratz et al.: Nucl.\ Phys.\ B 356, 332 (1991) 286
  \item M. L\"uscher: Commun.\ Math.\ Phys.\ 105, 153 (1986) 286, 287
  \item M. L\"uscher: Nucl.\ Phys.\ B 354, 531 (1991) 287
  \item M. L\"uscher: Nucl.\ Phys.\ B 364, 237 (1991) 287
  \item M. L\"uscher and U. Wolff: Nucl.\ Phys.\ B 339, 222 (1990) 288
  \item C. R. Gattringer and C. B. Lang: Nucl.\ Phys.\ B 391, 463 (1993) 288, 289
  \item J. Nishimura: Phys.\ Lett.\ B 294, 375 (1992) 289
  \item M. G\"ockeler, H. Kastrup, J. Westphalen, and F. Zimmermann: Nucl.\ Phys.\ B 425, 413 (1994) 289
  \item S. Aoki et al.: Phys.\ Rev.\ D 67, 014502 (2003) 289
  \item K. Rummukainen and S. Gottlieb: Nucl.\ Phys.\ B 450, 397 (1995) 289
  \item C. Kim, C. T. Sachrajda, and S. R. Sharpe: Nucl.\ Phys.\ B 727, 218 (2005) 289
  \item S. Aoki et al.: PoS LAT2006, 110 (2006) 289
  \item S. Aoki et al.: Phys.\ Rev.\ D 76, 094506 (2007) 289
  \item C. McNeile and C. Michael: Phys.\ Lett.\ B 556, 177 (2006) 289
  \item R. P. Feynman: \emph{Photon-Hadron Interactions} (Benjamin, New York 1972) 290
  \item D. Diakonov: Prog.\ Part.\ Nucl.\ Phys.\ 51, 173 (2003) 290
  \item X. Ji: Ann.\ Rev.\ Nucl.\ Part.\ Sci.\ 54, 413 (2004) 290
  \item K. G. Wilson: Phys.\ Rev.\ 179, 1499 (1969) 290, 295
  \item K. G. Wilson and W. Zimmermann: Commun.\ Math.\ Phys.\ 24, 87 (1972) 290, 295
  \item W. Zimmermann: Ann.\ Phys.\ 77, 570 (1973) 290, 295
  \item E. Witten: Nucl.\ Phys.\ B 120, 189 (1977) 290, 295
  \item J. Gounaris and J. J. Sakurai: Phys.\ Rev.\ Lett.\ 21, 244 (1968) 291
  \item M. F. Heyn and C. B. Lang: Z.\ Phys.\ C 7, 169 (1981) 291
  \item J. F. de Troc\'oniz and F. J. Yndur\'ain: Phys.\ Rev.\ D 65, 093001 (2002) 291
  \item H. Leutwyler: in \emph{Continuous Advances in QCD 2002}: Conference Proceedings Minneapolis, Minnesota, 2002, edited by K. A. Olive, M. A. Shifman, and M. B. Voloshin (World Scientific, Singapore 2002) 291, 294
  \item C. W. Bernard, T. Draper, G. Hockney, and A. Soni: in \emph{Lattice Gauge Theory: A Challenge in Large Scale Computing}, edited by B. Bunk, K. H. M\"utter, and K. Schilling (Plenum, New York 1986) 292
  \item G. W. Kilcup et al.: Phys.\ Lett.\ B 164, 347 (1985) 292
  \item D. Br\"ommel et al.: Eur.\ Phys.\ J.\ C 51, 335 (2007) 293, 294
  \item G. Martinelli and C. T. Sachrajda: Nucl.\ Phys.\ B 306, 865 (1988) 294
  \item T. Draper, R. Woloshyn, W. Wilcox, and K.-F. Liu: Nucl.\ Phys.\ B 318, 319 (1989) 294
  \item J. van der Heide, J. Koch, and E. Laermann: Phys.\ Rev.\ D 69, 094511 (2004) 294
  \item S. Capitani, C. Gattringer, and C. B. Lang: Phys.\ Rev.\ D 73, 034505 (2005) 294
  \item F. D. R. Bonnet et al.: Phys.\ Rev.\ D 72, 054506 (2005) 294
  \item S. Hashimoto et al.: PoS LAT2005, 336 (2005) 294
  \item H. Gausterer and C. B. Lang: Z.\ Phys.\ C 28, 475 (1985) 294
  \item G. Colangelo, J. Gasser, and H. Leutwyler: Nucl.\ Phys.\ B 603, 125 (2001) 294
  \item J. Bijnens: PoS LATTICE2007, 004 (2007) 294
  \item N. Cabibbo: Phys.\ Rev.\ Lett.\ 10, 531 (1963) 294
  \item M. Kobayashi and T. Maskawa: Prog.\ Theor.\ Phys.\ 49, 652 (1973) 294
  \item J. Buras: in \emph{Proceedings of the International School of Subnuclear Physics: 38th Course: Theory and Experiment Heading for New Physics}, Erice, Italy, 27 Aug--5 Sep 2000, edited by A. Zichichi (World Scientific, Singapore 2001) 295
  \item J. Buras: in \emph{Lectures given at the European CERN School}, Saint Feliu de Guixols, 2004 (CERN, Geneva, Switzerland 2005) 295
  \item Z. Ligeti: PoS LAT2005, 012 (2005) 295
  \item M. Okamoto: PoS LAT2005, 013 (2005) 295
  \item 5th International Workshop on the CKM Unitarity Triangle: \url{http://ckm2008.roma1.infn.it} (2008) 295
  \item A. Soni: Pramana 62, 415 (2004) 295
  \item S. Capitani and L. Giusti: Phys.\ Rev.\ D 64, 014506 (2001) 296
  \item L. Maiani, G. Martinelli, G. C. Rossi, and M. Testa: Phys.\ Lett.\ B 176, 445 (1986) 296
  \item L. Maiani, G. Martinelli, G. C. Rossi, and M. Testa: Nucl.\ Phys.\ B 289, 505 (1987) 296
  \item C. Bernard, T. Draper, G. Hockney, and A. Soni: Nucl.\ Phys.\ B (Proc.\ Suppl.) 4, 483 (1988) 296
  \item L. Maiani and M. Testa: Phys.\ Lett.\ B 245, 585 (1990) 296
  \item L. Lellouch and M. L\"uscher: Commun.\ Math.\ Phys.\ 219, 31 (2001) 296
  \item C.-J. Lin, G. Martinelli, C. Sachrajda, and M. Testa: Nucl.\ Phys.\ B 619, 467 (2001) 296
  \item L. Lellouch: PoS LATTICE2008, 015 (2009) 296
  \item T. DeGrand and C. DeTar: \emph{Lattice Methods for Quantum Chromodynamics} (World Scientific, Singapore 2006) 296
\end{enumerate}
